% Options for packages loaded elsewhere
\PassOptionsToPackage{unicode}{hyperref}
\PassOptionsToPackage{hyphens}{url}
%
\documentclass[
  10pt,
  a4paper,
]{article}
\usepackage{amsmath,amssymb}
\usepackage{iftex}
\ifPDFTeX
  \usepackage[T1]{fontenc}
  \usepackage[utf8]{inputenc}
  \usepackage{textcomp} % provide euro and other symbols
\else % if luatex or xetex
  \usepackage{unicode-math} % this also loads fontspec
  \defaultfontfeatures{Scale=MatchLowercase}
  \defaultfontfeatures[\rmfamily]{Ligatures=TeX,Scale=1}
\fi
\usepackage{lmodern}
\ifPDFTeX\else
  % xetex/luatex font selection
  \setmainfont[]{DejaVu Serif}
  \setmonofont[]{DejaVu Sans Mono}
  \setmathfont[]{Latin Modern Math}
\fi
% Use upquote if available, for straight quotes in verbatim environments
\IfFileExists{upquote.sty}{\usepackage{upquote}}{}
\IfFileExists{microtype.sty}{% use microtype if available
  \usepackage[]{microtype}
  \UseMicrotypeSet[protrusion]{basicmath} % disable protrusion for tt fonts
}{}
\makeatletter
\@ifundefined{KOMAClassName}{% if non-KOMA class
  \IfFileExists{parskip.sty}{%
    \usepackage{parskip}
  }{% else
    \setlength{\parindent}{0pt}
    \setlength{\parskip}{6pt plus 2pt minus 1pt}}
}{% if KOMA class
  \KOMAoptions{parskip=half}}
\makeatother
\usepackage{xcolor}
\usepackage[margin=20mm]{geometry}
\usepackage{longtable,booktabs,array}
\usepackage{calc} % for calculating minipage widths
% Correct order of tables after \paragraph or \subparagraph
\usepackage{etoolbox}
\makeatletter
\patchcmd\longtable{\par}{\if@noskipsec\mbox{}\fi\par}{}{}
\makeatother
% Allow footnotes in longtable head/foot
\IfFileExists{footnotehyper.sty}{\usepackage{footnotehyper}}{\usepackage{footnote}}
\makesavenoteenv{longtable}
\setlength{\emergencystretch}{3em} % prevent overfull lines
\providecommand{\tightlist}{%
  \setlength{\itemsep}{0pt}\setlength{\parskip}{0pt}}
\setcounter{secnumdepth}{-\maxdimen} % remove section numbering
\usepackage{amsmath,amssymb,amsthm,mathtools,graphicx,booktabs,longtable,needspace}
\usepackage{fvextra}
\DefineVerbatimEnvironment{Highlighting}{Verbatim}{commandchars=\\\{\},breaklines=true,breakanywhere=true,fontsize=\small}
\setlength{\emergencystretch}{3em}
\allowdisplaybreaks
\raggedbottom\sloppy
\widowpenalty=10000\clubpenalty=10000
\providecommand{\Tr}{\operatorname{Tr}}
\providecommand{\tr}{\operatorname{Tr}}
\ifLuaTeX
  \usepackage{selnolig}  % disable illegal ligatures
\fi
\usepackage{bookmark}
\IfFileExists{xurl.sty}{\usepackage{xurl}}{} % add URL line breaks if available
\urlstyle{same}
\hypersetup{
  pdftitle={Angular contrast and rounding},
  pdfauthor={Kirtus Justus},
  hidelinks,
  pdfcreator={LaTeX via pandoc}}

\title{Angular contrast and rounding}
\author{Kirtus Justus}
\date{}

\begin{document}
\maketitle

\subsection{1. Quantum model, notation and sufficient two-dimensional reduction}\label{quantum-model-notation-and-sufficient-two-dimensional-reduction}

Let \texttt{P=\textbar{}p\textgreater{}\textless{}p\textbar{}} be a rank-one projector and \texttt{R=I−P} on the specified joint-probe Hilbert space. Let \texttt{ρ} be a normalized state and \texttt{0≤E≤I} a \emph{common} binary effect, with

\[
 s=\mathrm{Tr}(P\rho)=\frac{131}{320},\qquad h=\mathrm{Tr}(EP)=\frac{99}{100},\qquad
 q=q(0)=\frac{81711338569043}{224995905216000}.
\]

For the phase family \texttt{U\_φ=P+e\^{}\{iφ\}R}, write \texttt{q(φ)=Tr(E\ U\_φ\ ρ\ U\_φ†)}. Put

\[
 a=sh,\quad d=s(1-h),\quad t=a-q>0,\quad
 b=\mathrm{Tr}(E R\rho R),\quad c=\mathrm{Tr}(E P\rho R),\quad z=2\operatorname{Im}c.
\]

The ordinary readout fixes \texttt{Re\ c=−(t+b)/2}. Positivity of \texttt{E} and \texttt{I−E} gives

\[
 |c|^2\le\min\{ab,\ d(1-s-b)\},\qquad 0\le b\le 1-s.
\]

These conditions are \textbf{sufficient}, not just necessary, because every admissible pair \texttt{(b,c)} is realized by the pure state \texttt{\textbar{}ψ\textgreater{}=√s\textbar{}p\textgreater{}+√(1−s)\textbar{}u\textgreater{}} and the two-dimensional effect

\[
 E_{b,c}=\begin{pmatrix}
 h & \overline c/\sqrt{s(1-s)}\\
 c/\sqrt{s(1-s)} & b/(1-s)
 \end{pmatrix},
\]

embedded in the joint-probe space. Both determinant conditions for \texttt{E} and \texttt{I−E} are precisely the inequalities above. Their joint feasible set is convex.

Equivalently, the allowed \texttt{(b,z)} are the intersection of two disks:

\[
 (b-C_1)^2+z^2\le R_1^2,\quad (b-C_2)^2+z^2\le R_2^2,
\] \[
 C_1=a+q,\ R_1=2\sqrt{aq},\quad
 C_2=-(t+2d),\ R_2=2\sqrt{d(1-q)}.
\]

They meet at \texttt{b₀=(1−s)(1−h)} and \texttt{z=±z₀}, with

\[
 z_0=\sqrt{4ab_0-(t+b_0)^2},\qquad
 \kappa=\frac{z_0}{t+b_0}=1.77548900478028319584\ldots.
\]

\subsection{\texorpdfstring{2. Sharp \textbf{two-sided} contrast for a deterministic coherent offset}{2. Sharp two-sided contrast for a deterministic coherent offset}}\label{sharp-two-sided-contrast-for-a-deterministic-coherent-offset}

Let the intended relative phase be \texttt{π}; assume that its \textbf{only} error is a coherent offset \texttt{±δ} in the very same family \texttt{U\_φ}, with the \textbf{same} \texttt{ρ}, \texttt{E} and \texttt{P} in both arms. Then, upon choosing the adverse sign of the imaginary interference,

\[
 \Delta(\delta):=q(\pi\pm\delta)-q(0)
 =A(t+b)-Bz,\qquad A=1+\cos\delta,\quad B=\sin\delta
\]

for \texttt{0≤δ≤π} (the signs are interchangeable because \texttt{z} can take either sign). Set \texttt{H=√(A²+B²)=2cos(δ/2)}.

\textbf{Sharp minimum.} Let

\[
 \delta_1=2\arccos[(C_1-b_0)/R_1]=12.7591826805379687\ldots^\circ.
\]

For \texttt{0≤δ≤δ₁}, the supporting point of the first disk is feasible for both, giving

\[
 \Delta_{\min}(\delta)=A(t+C_1)-R_1H
 =4a\cos^2(\delta/2)-4\sqrt{aq}\cos(\delta/2).
\]

For \texttt{δ₁≤δ≤δ\_c} the supporting point would leave the disk intersection; the minimum is at its shared upper point \texttt{(b₀,z₀)}:

\[
 \Delta_{\min}(\delta)=A(t+b_0)-Bz_0.
\]

Its zero occurs sharply at

\[
 \delta_c=2\arctan(1/\kappa)=58.77860795856966018\ldots^\circ.
\]

Every compatible state/effect yields strictly positive contrast for \texttt{\textbar{}δ\textbar{}\textless{}δ\_c}. At \texttt{δ\_c} an admissible two-dimensional \textbf{projective} effect gives zero, so this is an open condition.

\textbf{New sharp upper endpoint.} On the same full interval \texttt{0≤δ≤δ\_c}, maximizing the contrast over the \emph{second} disk gives the simultaneously attainable bound

\[
 \boxed{\displaystyle
 \Delta_{\max}(\delta)=A(t+C_2)+R_2H
 =4\sqrt{d(1-q)}\cos(\delta/2)-4d\cos^2(\delta/2).}
\]

Its supporting point is

\[
 (b,z)=\left(C_2+R_2\frac A H,\ -R_2\frac B H\right).
\]

This point belongs to the first disk whenever \texttt{b≥b₀}. For the supplied inputs it remains feasible through \texttt{δ=113.20968269065858...°}, well beyond \texttt{δ\_c}. The upper formula therefore holds throughout the guaranteed-positive interval. The constructed two-dimensional \texttt{E\_\{b,c\}} has \texttt{0≤E≤I} and attains the rate. The image of the convex disk intersection under the linear contrast is a full interval, so \textbf{every intermediate contrast is attainable} as well.

\begin{longtable}[]{@{}
  >{\raggedleft\arraybackslash}p{(\columnwidth - 4\tabcolsep) * \real{0.3333}}
  >{\raggedleft\arraybackslash}p{(\columnwidth - 4\tabcolsep) * \real{0.3333}}
  >{\raggedleft\arraybackslash}p{(\columnwidth - 4\tabcolsep) * \real{0.3333}}@{}}
\toprule\noalign{}
\begin{minipage}[b]{\linewidth}\raggedleft
Offset
\end{minipage} & \begin{minipage}[b]{\linewidth}\raggedleft
Sharp lower increase (percentage points)
\end{minipage} & \begin{minipage}[b]{\linewidth}\raggedleft
Sharp upper increase (percentage points)
\end{minipage} \\
\midrule\noalign{}
\endhead
\bottomrule\noalign{}
\endlastfoot
0° & 8.653592287309 & 18.786135457214 \\
10° & 8.006123191762 & 18.720856010456 \\
20° & 6.398293220245 & 18.525231209691 \\
30° & 4.697644815950 & 18.199908655488 \\
45° & 2.168778595748 & 17.471285101295 \\
55° & 0.572299788190 & 16.827620196169 \\
58° & 0.116299082295 & 16.610292664719 \\
\end{longtable}

These are percentages of theoretical binary \emph{probabilities}, not percentages measured on devices. The decimal entries in this table are rounded descriptions, not lower/upper decimal certificates. The companion verifier emits \textbf{downward-rounded lower certificates} for \texttt{Δ\_min} separately.

\subsection{3. Sharp calibration criterion for an independently varying phase offset}\label{sharp-calibration-criterion-for-an-independently-varying-phase-offset}

Suppose one fixed pre-phase conditional \texttt{ρ} and one fixed \texttt{E} apply to all trials in this comparison, and the distribution of actual offsets \texttt{ε} is transferable to that same response independently of hidden state/effect changes. Define the \emph{population} moments

\[
 C=\mathbb E\cos\varepsilon,\qquad S=\mathbb E\sin\varepsilon,
\qquad C^2+S^2\le1.
\]

No values of these moments are assumed known from experiment. Averaging the fixed response gives exactly

\[
 \overline\Delta=(1+C)(t+b)-S z.
\]

Put \texttt{A=1+C}, \texttt{B=\textbar{}S\textbar{}}, \texttt{H=√(A²+B²)} and \texttt{b\_*\ =\ C₁−R₁\ A/H} when \texttt{H\textgreater{}0}. The \textbf{sharp minimum} over all state/effect pairs compatible with the three supplied population inputs is

\[
 \boxed{\displaystyle
 M(C,S)=\begin{cases}
 A(t+C_1)-R_1H, & b_*\le b_0,\\
 A(t+b_0)-Bz_0, & b_*>b_0.
 \end{cases}}
\]

If \texttt{A=B=0} (the phase is reversed relative to the intended reflection with probability one), take \texttt{M=0}. The proof is the same supporting-line optimization over the two-disk intersection, and both branches are attainable.

The exact \textbf{strict sign criterion} is

\[
 \boxed{1+C>\kappa|S|.}
\]

Necessity as well as sufficiency follows because \texttt{max\ z/(t+b)=κ} at the disk intersection, and \texttt{t+b\textgreater{}0} throughout the feasible set. In the symmetric case \texttt{S=0} the sharp lower contrast simplifies to

\[
 \boxed{M(C,0)=\frac{1+C}{2}\;4(a-\sqrt{aq}).}
\]

Simultaneous, independently justified confidence bounds \texttt{C≥C\_L} and \texttt{\textbar{}S\textbar{}≤S\_U} therefore give the sufficient physical sign criterion \texttt{1+C\_L\textgreater{}κ\ S\_U} \emph{only when} the common-state/effect and transfer hypotheses have been independently established. The rectangular bound can be conservative if it includes moment pairs that no phase distribution can realize.

\textbf{Nontransfer control.} Two hidden conditional subpopulations can have the same \texttt{s,h,q} but opposite imaginary interference coefficients. If positive and negative angle errors are systematically correlated with these subpopulations, their pooled error histogram can have \texttt{S=0} while both arms have the adverse response. At the sign boundary their average contrast can be zero. Thus an unlabelled marginal angle histogram is \textbf{not} evidence for the independent-transfer condition, and may not be inserted into \texttt{M(C,S)} without justification.

\end{document}
